\documentclass[10pt,a4paper]{amsart}
\usepackage{amsmath,amssymb,amsthm,graphicx}
\usepackage{verbatim}
\newenvironment{reference}{\comment}{\endcomment}
%\newenvironment{reference}{}{}
\newenvironment{slogan}{\comment}{\endcomment}
\newenvironment{history}{\comment}{\endcomment}

% For commutative diagrams we use Xy-pic
\usepackage[all]{xy}

% We use 2cell for 2-commutative diagrams.
\xyoption{2cell}
% Theorem environments.
%
\theoremstyle{plain}
\newtheorem{theorem}[subsection]{Theorem}
\newtheorem{proposition}[subsection]{Proposition}
\newtheorem{lemma}[subsection]{Lemma}

\theoremstyle{definition}
\newtheorem{definition}[subsection]{Definition}
\newtheorem{example}[subsection]{Example}
\newtheorem{exercise}[subsection]{Exercise}
\newtheorem{situation}[subsection]{Situation}

\theoremstyle{remark}
\newtheorem{remark}[subsection]{Remark}
\newtheorem{remarks}[subsection]{Remarks}

\numberwithin{equation}{subsection}

% Macros
%
\def\lim{\mathop{\mathrm{lim}}\nolimits}
\def\colim{\mathop{\mathrm{colim}}\nolimits}
\def\Spec{\mathop{\mathrm{Spec}}}
\def\Hom{\mathop{\mathrm{Hom}}\nolimits}
\def\Ext{\mathop{\mathrm{Ext}}\nolimits}
\def\SheafHom{\mathop{\mathcal{H}\!\mathit{om}}\nolimits}
\def\SheafExt{\mathop{\mathcal{E}\!\mathit{xt}}\nolimits}
\def\Sch{\mathit{Sch}}
\def\Mor{\mathop{\mathrm{Mor}}\nolimits}
\def\Ob{\mathop{\mathrm{Ob}}\nolimits}
\def\Sh{\mathop{\mathit{Sh}}\nolimits}
\def\NL{\mathop{N\!L}\nolimits}
\def\CH{\mathop{\mathrm{CH}}\nolimits}
\def\proetale{{pro\text{-}\acute{e}tale}}
\def\etale{{\acute{e}tale}}
\def\QCoh{\mathit{QCoh}}
\def\Ker{\mathop{\mathrm{Ker}}}
\def\Im{\mathop{\mathrm{Im}}}
\def\Coker{\mathop{\mathrm{Coker}}}
\def\Coim{\mathop{\mathrm{Coim}}}

% Boxtimes
%
\DeclareMathSymbol{\boxtimes}{\mathbin}{AMSa}{"02}

%
% Macros for moduli stacks/spaces
%
\def\QCohstack{\mathcal{QC}\!\mathit{oh}}
\def\Cohstack{\mathcal{C}\!\mathit{oh}}
\def\Spacesstack{\mathcal{S}\!\mathit{paces}}
\def\Quotfunctor{\mathrm{Quot}}
\def\Hilbfunctor{\mathrm{Hilb}}
\def\Curvesstack{\mathcal{C}\!\mathit{urves}}
\def\Polarizedstack{\mathcal{P}\!\mathit{olarized}}
\def\Complexesstack{\mathcal{C}\!\mathit{omplexes}}
% \Pic is the operator that assigns to X its picard group, usage \Pic(X)
% \Picardstack_{X/B} denotes the Picard stack of X over B
% \Picardfunctor_{X/B} denotes the Picard functor of X over B
\def\Pic{\mathop{\mathrm{Pic}}\nolimits}
\def\Picardstack{\mathcal{P}\!\mathit{ic}}
\def\Picardfunctor{\mathrm{Pic}}
\def\Deformationcategory{\mathcal{D}\!\mathit{ef}}
\usepackage{cleveref}
\crefname{originallemma}{original lemma}{original lemmas}
\Crefname{originallemma}{Original lemma}{Original lemmas}
\makeatletter
\newlabel{divisors-lemma-extend-rational-map-after-modification}{{0C4V}{source}{}{sourcepointer-1}{}}
\newlabel{divisors-lemma-extend-rational-map-after-modification@cref}{{[originallemma][0][]0C4V}{[1][0][]source}{}{}{}}
\makeatother
\crefname{originalreference}{original reference}{original references}
\Crefname{originalreference}{Original reference}{Original references}
\makeatletter
\newlabel{lemma-dominate-by-normalized-blowing-up}{{0BBT}{source}{}{sourcepointer-2}{}}
\newlabel{lemma-dominate-by-normalized-blowing-up@cref}{{[originalreference][0][]0BBT}{[1][0][]source}{}{}{}}
\makeatother
\title{Extending a rational map by normalized blowups for a Nagata surface}
\author{The Stacks Project Authors}
\date{stacks-project:resolve.tex:local-snapshot-20261012; GNU FDL 1.2 or later}
\begin{document}
\maketitle
\noindent\textit{Editorial note.} All delivered mathematical text is the author's own, copied from the pinned Stacks Project chapter file resolve.tex; no statement or proof step was rewritten, reordered or removed.\par\medskip
\noindent\textit{Editorial note.} The counted proof is the author's complete argument: the proper modification over which the rational map is defined, and its domination by a normalization followed by normalized blowups in closed points.\par\medskip
\noindent\textit{Editorial note.} The label lemma-extend-rational-map-normalized-blowing-up names the counted statement, and the proof environment immediately following it in the source is delivered verbatim as its proof.\par\medskip
\noindent\textit{Editorial note.} The counted statement is carried by the author's own lemma environment in the chapter source.\par\medskip
\noindent\textit{Editorial note.} Pre-open gate: python operations/agent-authoring/tools/precard\_gate.py --source resolve.tex --label lemma-extend-rational-map-normalized-blowing-up --start 1389 --end 1411 --local-id 28122 returned RESULT=OK (L28122 free, occupancy FREE, public interval clear). Pre-assembly tag self-check: canonical tag resolved and not present in reference/public-items.jsonl.\par\medskip
\noindent\textit{Editorial note.} The statement and proof are copied from the pinned chapter file; the cited inputs are kept as external pointers with their canonical Stacks tags.\par\medskip
\noindent\textit{Allowed standard prerequisites (editorial).}\begin{itemize}
\item Extending a rational map after a modification. Divisors, the lemma producing a proper modification over which a rational map into a proper target is defined. Produces the proper morphism to the source which is an isomorphism over the domain of the rational map and carries a morphism agreeing with it.
\item Dominating a proper modification by normalized blowups. Same chapter, the lemma dominating a proper modification of a Nagata surface by a sequence of normalized blowups in closed points. Turns the modification into a normalization followed by a sequence of normalized blowups and provides the morphism to the target.
\end{itemize}\medskip
\ifdefined\hypertarget\hypertarget{sourcepointer-1}{}\fi
\noindent\textit{Source reference.} Original lemma 0C4V: sources/stacks-chow/tags.txt, source lines 15240--15240. Prerequisite pointer invoked at proof line(s) 1405; it names a standard fact of a sibling chapter that the counted argument uses, and every step of the counted proof built on it is retained verbatim in the delivered body.\par\medskip
\ifdefined\hypertarget\hypertarget{sourcepointer-2}{}\fi
\noindent\textit{Source reference.} Original reference 0BBT: sources/stacks-resolve/pinned-original-source.tex, source lines 1277--1277. Prerequisite pointer invoked at proof line(s) 1408; it names a standard fact of the same chapter that the counted argument uses, and every step of the counted proof built on it is retained verbatim in the delivered body.\par\medskip
\begin{lemma}
\label{lemma-extend-rational-map-normalized-blowing-up}
Let $S$ be a scheme. Let $X$ be a scheme over $S$ which is
Noetherian, Nagata, and has dimension $2$. Let $Y$ be a proper
scheme over $S$. Given an $S$-rational map $f : U \to Y$ from
$X$ to $Y$ there exists a sequence
$$
X_n \to X_{n - 1} \to \ldots \to X_1 \to X_0 \to X
$$
and an $S$-morphism $f_n : X_n \to Y$ such that $X_0 \to X$ is the
normalization, $X_{i + 1} \to X_i$ is the normalized blowing up of
$X_i$ at a closed point, and $f_n$ and $f$ agree.
\end{lemma}

\begin{proof}
Applying
Divisors, Lemma \ref{divisors-lemma-extend-rational-map-after-modification}
we find a proper morphism $p : X' \to X$ which is an isomorphism
over $U$ and a morphism $f' : X' \to Y$ agreeing with $f$ over $U$.
Apply Lemma \ref{lemma-dominate-by-normalized-blowing-up}
to the morphism $p : X' \to X$. The composition $X_n \to X' \to Y$ is
the desired morphism.
\end{proof}

\end{document}
